\section{Introduction}

Agent memory is not automatically useful. A fresh coding agent may inspect the
current repository, run focused tests, and reconstruct the necessary behavior
without predecessor context. In that setting, injecting older summaries can add
latency, consume context, or promote a stale assumption. A credible evaluation
must therefore preserve the ordinary capabilities of a no-memory agent instead
of weakening it to manufacture a benefit.

\memorix{} is a project memory system for coding agents. Its intended role is
not to replay every prior interaction. It forms a task workset that combines
current code state with bounded durable project evidence, warnings for suspect
claims, and optional knowledge and workflow context. Current source has
precedence over narrative memory.

This paper introduces the evaluation design for that claim. The central question
is conditional: for a declared cross-session dependency after controlled project
evolution, does freshness-aware project memory improve the next agent's
engineering outcome relative to a capable fresh agent with the same current
repository access? Equally important, does optional use abstain when current
source is already enough?

Our contributions are an evaluation protocol for a narrow fresh-agent transfer
question, a strong no-memory control, a separation between canonical retrieval
and native product behavior, a 144-run public reproducible cohort plus a
separately frozen 72-run cross-model replication, and an evidence boundary
that rejects local diagnostics as confirmatory results. We do not claim that
software evolution, dynamic memory, or sequential coding tasks are new in
themselves; adjacent benchmarks already study those broader settings.
